\documentclass[11pt]{article}
\usepackage[margin=1in]{geometry}
\usepackage{graphicx}
\usepackage{amsmath,amssymb}
\usepackage{booktabs}
\usepackage{xcolor}
\usepackage{hyperref}
\hypersetup{colorlinks=true,linkcolor=blue,urlcolor=blue}
\usepackage{caption}
\captionsetup{font=small}

\newcommand{\op}[1]{\texttt{#1}}
\graphicspath{{../archive/}}

\title{\textbf{Active Matter $\times$ MPM: Cells in a Deformable Blob}\\[2pt]
\large A Plexus prototype toward embryogenesis / morphogenesis}
\author{Plexus prototype \texttt{prototype/embryogenesis}}
\date{\today}

\begin{document}
\maketitle

\begin{abstract}
We couple \emph{active-matter agents} (self-propelled, aligning ``cells'') to a \emph{Material
Point Method} (MLS-MPM) continuum ``blob'' of water or elastic tissue. The two subsystems talk
through the shared MPM background grid: the material \emph{drags} the cells and its
surface-tension interface \emph{confines} them; the cells scatter momentum back and
\emph{deform} the material. The blob rotates slowly. Everything is built from Plexus operators;
the only new physics is three small operators (\op{mpm\_to\_agent}, \op{agent\_to\_mpm},
\op{mpm\_spin}). We show stable 2D prototypes for a water disc, an elastic disc, and a
four-cell-type showcase, and document the confinement (``escape'') failure and its fix. The
system is the substrate for an agentic exploration loop that maps its morphogenetic regimes.
\end{abstract}

\section{Goal}
Multi-type cells migrate collectively inside a large disc of MPM material (water or elastic).
The material drags the cells; the cells create stress/deformation in the material; the cells
cannot leave the disc, which is held together by surface tension. Ultimately we want to
reproduce embryogenesis / morphogenesis dynamics. This note records the working prototypes and
the code they required.

\section{Model}
Two node sets share one world and one clock.

\paragraph{Cells (active matter).} A set \op{agent} of self-propelled particles carrying a unit
heading $\mathbf n_i$. Motion is heading-kinematic (overdamped, first order): the heading is
steered by polar alignment (Vicsek), chemotaxis, and angular noise, while \op{glide} moves each
cell along its heading at speed $v_0^{(\tau_i)}$ (per type $\tau_i$) and \op{repel} enforces a
hard core:
\begin{equation}
\dot{\mathbf r}_i = v_0^{(\tau_i)}\,\mathbf n_i + \sum_j \mathbf f_{ij}, \qquad
\dot{\mathbf n}_i = \Gamma\!\!\sum_{j\in\partial i}\!\frac{\mathbf n_j-(\mathbf n_j\!\cdot\!\mathbf n_i)\mathbf n_i}{r_{ij}}
+ \omega\,\widehat{\nabla c} + \sqrt{2D_r}\,\xi_i .
\end{equation}

\paragraph{Material (MLS-MPM).} A disc of material points (set \op{mpm\_particle}, child of a
\op{cell}) with deformation gradient $F$ and affine velocity $C$. The decomposed solver is the
standard cycle \op{mpm\_strain} $\to$ \op{p2g} $\to$ \op{mpm\_grid\_update} $\to$ \op{g2p}, run
for $\lfloor \Delta t/\Delta t_\mathrm{sub}\rfloor$ substeps per frame. Water is a liquid band
($\mu=0$, pressure only); elastic tissue is fixed-corotated with a thin liquid \emph{skin}. The
liquid colour field $c$ drives the continuum surface force (CSF) that holds the blob's boundary.

\paragraph{Two-way coupling (new).} Both directions route through the MPM grid velocity
$\mathbf v_g$ and mass/momentum $m_g,\,\mathbf p_g$:
\begin{align}
\text{material}\to\text{cell:}\quad
&\dot{\mathbf r}_i \mathrel{+}= k\,\mathbf v_g(\mathbf r_i) \;+\; \kappa\,\nabla c(\mathbf r_i)
&&(\op{mpm\_to\_agent}) \label{eq:m2a}\\
\text{cell}\to\text{material:}\quad
&m_g \mathrel{+}= \textstyle\sum_i w\,\mu_a,\quad
\mathbf p_g \mathrel{+}= \textstyle\sum_i w\,\mu_a v_0\mathbf n_i
&&(\op{agent\_to\_mpm}) \label{eq:a2m}
\end{align}
using the \emph{same} quadratic B-spline weights $w$ MPM uses for its own transfers. In
\eqref{eq:m2a} the first term is fluid drag (the cell is advected by the flow, $k{=}1$ = fully
carried) and the second is the \textbf{surface-tension confinement}: $c$ is $\sim\!1$ inside the
fluid and $\sim\!0$ outside, so $\nabla c$ points \emph{inward} at the interface and pushes cells
back in --- a soft membrane, not a hard wall. In \eqref{eq:a2m} cells have no physical mass, so an
effective mass $\mu_a$ (times a gain $k$) sets the push; scattering onto $m_g,\mathbf p_g$
\emph{after} \op{p2g} but \emph{before} \op{mpm\_grid\_update} makes the grid solve on the
combined (material $+$ cell) momentum, so \op{g2p} carries the cells' push into real deformation.
Rotation is \op{mpm\_spin}: a body force driving the disc toward solid-body rotation
$\mathbf v_\mathrm{rot}=\omega\,\hat{\mathbf z}\times(\mathbf r-\mathbf c)$.

\paragraph{One clock.} The engine integrates on a single global $\Delta t$; we make the MPM clock
master ($\Delta t=2\!\times\!10^{-3}$, $\Delta t_\mathrm{sub}=2\!\times\!10^{-4}$) and express the
cell gains in that time unit.

\section{Results}
All panels: cells are coloured dots by type, the MPM material is the translucent blue blob,
black background. Time runs left$\to$right.

\begin{figure}[h]\centering
\includegraphics[width=\linewidth]{agent_mpm_disc_water_v3/blob_evolution.png}
\caption{\textbf{Water disc.} Two cell types swim in a cohesive rotating water blob, form
aligned streams, and gently lobe the boundary. $0\%$ escape.}
\end{figure}

\begin{figure}[h]\centering
\includegraphics[width=\linewidth]{agent_mpm_disc_elastic_v3/blob_evolution.png}
\caption{\textbf{Elastic disc.} A deformable tissue with a liquid skin holds a rounder shape; a
coherent stream of cells (polar order $\approx0.46$) migrates across it. $0.25\%$ escape.}
\end{figure}

\begin{figure}[h]\centering
\includegraphics[width=\linewidth]{agent_mpm_disc_4types_show/blob_evolution.png}
\caption{\textbf{Four-type showcase} ($4000$ cells). Per-type behaviour via masked operator
instances: a\,(red) flocker, b\,(gold) disperser, c\,(green) chemotactic aggregator, d\,(white)
noisy. All confined and carried by the rotating water blob.}
\end{figure}

\section{The escape problem and its fix}
The first attempt let nearly all cells leak out and the water disc foam apart
(Fig.~\ref{fig:fail}). Two coupled causes: (i) an under-resolved liquid (too few particles per
grid cell) with \emph{zero} viscous drag, agitated by spin and cell push, fragments; (ii) once
the interface is ragged and the confinement gain is weak, cells walk straight through it. The fix
is a combination, summarised in Table~\ref{tab:tune}: fewer grid cells (more particles/cell),
viscous drag, stronger surface tension, gentler spin/push, and a much stronger confinement gain.

\begin{figure}[h]\centering
\includegraphics[width=0.85\linewidth]{water_v1_ESCAPE_HIST/fig_mpm_particle_evolution.png}
\caption{\textbf{Failure kept on purpose (v1).} The water disc foams, fragments and spreads;
cells escape to the walls.}\label{fig:fail}
\end{figure}

\begin{table}[h]\centering
\begin{tabular}{lccccl}
\toprule
run & escaped & disc growth & aniso & polar & key change \\
\midrule
water v1 (fail) & $\sim0.95$ & $0.30$ & $0.60$ & --- & baseline: n\_grid 96, drag 0, confine 30 \\
water v2 & $0.025$ & $0.042$ & $0.14$ & $0.14$ & n\_grid 64, drag 0.25, ST 120, confine 70 \\
water v3 & $\mathbf{0.000}$ & $0.057$ & $0.14$ & $0.14$ & drag 0.4, ST 160--300, confine 140 \\
elastic v3 & $0.003$ & $0.049$ & $0.08$ & $\mathbf{0.46}$ & elastic core $+$ liquid skin \\
\bottomrule
\end{tabular}
\caption{Confinement / stability sweep. ``escaped'' = fraction of cells outside the blob at the
end; ``aniso'' = blob boundary roughness (0 = perfect disc).}\label{tab:tune}
\end{table}

\section{Compute budget}
On one RTX~A6000, the four-type showcase ($4000$ cells, $16{,}000$ material points, $n_\mathrm{grid}=64$)
simulates $1500$ frames in $\approx100$\,s. This leaves ample headroom for the $\le 20$\,min-per-job
L4 budget of the agentic exploration loop.

\section{Next: agentic exploration}
The prototype is the substrate for an autonomous loop (a clone of \op{prototype/active\_matter2})
in which an agent proposes parameter variations, runs $8$ jobs in parallel on L4, and maps the
morphogenetic regimes (mixing vs.\ segregation, blob rounding vs.\ lobing, engulfment,
rotation-shear banding), with a batch montage as the visual feedback. In 3D the same operators
drive a 2$\times$2 summary movie: cell trajectories with tails, a transparent deformation volume,
the affine-velocity field, and a combined view.

\end{document}
